\chapter{Conclusion and Future Work}

\section{Summary}

This project built an emotion detection system for social media text that does something most
existing systems do not: it knows when not to answer. Instead of always returning a confident
label, the system can output Ambiguous when the input text is genuinely unclear or sits in a
boundary region between emotion classes.

The system combines three types of features. Semantic embeddings from the
\texttt{all-mpnet-base-v2} model capture the meaning of the text. TF-IDF features capture the
specific words and phrases present. Topological features from Vietoris-Rips persistent homology
capture the geometric structure of the neighbourhood around the input in the embedding space.

Three XGBoost classifiers trained in a one-vs-rest configuration produce probabilities for the
Positive, Negative, and Neutral classes. On top of this, an ambiguity detection layer combines
two signals: the number of H1 loops in the local persistence diagram and the confidence gap
between the top two class probabilities. Three decision rules convert these signals into a final
output, which can be Positive, Negative, Neutral, or Ambiguous.

\section{What Was Achieved}

The key achievements of this project are:

A working emotion classifier was built on the full GoEmotions dataset (43,410 training samples,
5,427 test samples), achieving an accuracy of 0.59 and a weighted F1 score of 0.60. These
numbers reflect the genuine difficulty of the dataset, not a fundamental failure of the approach.

A topological ambiguity detection mechanism was introduced. This mechanism uses local
Vietoris-Rips persistent homology computed fresh at inference time for every input. The number
of H1 loops in the persistence diagram provides a geometric signal of ambiguity that is
independent of the classifier's output.

The system correctly identifies genuinely ambiguous sentences. ``i dont know how i feel''
returns \textsc{ambiguous} due to an extreme loop count of 15. ``i will kill you'' returns
\textsc{ambiguous} due to a very low confidence gap of 0.06. ``do this'' returns
\textsc{ambiguous} due to moderate topology combined with low classifier confidence.

The system provides interpretable outputs. When it returns \textsc{ambiguous}, it explains
which rule triggered the decision. This is important for practical deployment where users need
to understand and trust the system's decisions.

\section{Advantages of the Approach}

The main advantage of using topology is that it provides a signal that is independent of the
classifier. A classifier can be trained to be confident on ambiguous inputs if it has seen many
similar examples. Topology looks at the structure of the data itself.

The approach is also principled. The connection between geometric complexity (H1 loops) and
semantic ambiguity (text near the boundary of emotion clusters) is not arbitrary. It reflects a
real property of the embedding space: texts that are semantically ambiguous are found in
geometrically complicated regions.

The system is also explainable. Every output comes with a reason. This is a property that most
deep learning systems do not have.

\section{Limitations}

The system has the following known limitations:

\begin{itemize}
    \item \textbf{Threshold Calibration:} The thresholds for ambiguity detection (14, 10, 0.30,
    0.10) were set manually based on experimentation. Without a labelled ambiguity dataset,
    these cannot be calibrated systematically.

    \item \textbf{Social Ambiguity Blind Spots:} The system cannot detect social ambiguity
    that is not visible in the embedding space. ``I am fine'' is socially complex but does
    not appear geometrically ambiguous because the training data biases its neighbourhood
    toward clearly Positive samples.

    \item \textbf{Slang and Informal Language:} The system struggles with slang and informal
    abbreviations. Words like ``yolo'', ``fomo'', ``lol'', and ``idk'' carry emotional meaning
    but may not be well-represented in the training data. The embedding model may not
    place such terms accurately, leading to incorrect or ambiguous predictions.

    \item \textbf{Inference Bottleneck:} Inference is slower than a standard classifier because
    of the live topology computation. For large-scale systems, this would need to be
    optimized.

    \item \textbf{Language Constraint:} The system works only on English text.
\end{itemize}

\section{Future Work}

There are several directions where this work can be extended:

\begin{itemize}
    \item \textbf{Labelled Ambiguity Dataset:} The most important improvement would be
    creating a labelled ambiguity dataset. If a set of social media posts were annotated by
    humans for whether they are genuinely ambiguous, the ambiguity thresholds could be
    calibrated properly and the system's performance on ambiguity detection could be
    measured numerically.

    \item \textbf{Faster Topology Computation:} A faster topology computation method would
    make the system more practical for deployment. One option is to precompute topological
    signatures for regions of the embedding space and do a lookup at inference time instead
    of computing fresh persistence for every input.

    \item \textbf{Fine-Grained Emotion Classes:} Extending the system to handle fine-grained
    emotion classes, rather than the collapsed three-class version, would be a more
    challenging but more useful task. The GoEmotions dataset supports 28 classes, and a
    system that can detect ambiguity at this level would be much more informative.

    \item \textbf{Sarcasm and Irony Detection:} Adding a dedicated sarcasm and irony
    detection module would address many of the cases where social ambiguity is not
    captured by the topological signal. Sarcasm in particular is common in Reddit comments
    and is a major source of error in sentiment analysis systems.

    \item \textbf{Multilingual Support:} Supporting multiple languages through multilingual
    embeddings would make the system applicable to a much wider range of social media
    content.

    \item \textbf{User Interface:} Finally, a web interface would make the system accessible
    to non-technical users. The current command-line interface is sufficient for development
    and testing but not for general use.
\end{itemize}

\section{Final Thoughts}

Emotion detection is a genuinely hard problem. Human language is full of nuance,
context-dependence, and social complexity. A system that always gives a confident answer is
not always being honest.

This project shows that it is possible to build a system that knows what it does not know and
says so clearly. The combination of geometric structure from topology and probabilistic output
from a classifier is a promising direction for building more reliable and trustworthy natural
language processing systems. The GoEmotions dataset, with all its noise and imbalance,
provides a realistic and challenging test bed for exploring this direction.

The work here is a starting point. There is much more to be done in characterizing ambiguity in
text, calibrating geometric signals, and connecting mathematical structures from topology to the
practical needs of emotion understanding.
